% Generated draft: requires mathematical and source review.
\documentclass[11pt]{article}
\usepackage[margin=25mm]{geometry}
\usepackage{fontspec}
\setmainfont{Latin Modern Roman}
\usepackage{amsmath,amssymb,mathtools,unicode-math}
\setmathfont{Latin Modern Math}
\usepackage{enumitem,needspace,xurl,hyperref}
\setlength{\parindent}{0pt}
\setlength{\parskip}{5pt}
\setlength{\emergencystretch}{3em}
\newcommand{\N}{\mathbb{N}}
\newcommand{\Z}{\mathbb{Z}}
\newcommand{\Q}{\mathbb{Q}}
\newcommand{\R}{\mathbb{R}}
\newcommand{\C}{\mathbb{C}}
\newcommand{\F}{\mathbb{F}}
\newcommand{\A}{\mathbb{A}}
\newcommand{\PP}{\mathbb{P}}
\newcommand{\E}{\mathbb{E}}
\newcommand{\eps}{\varepsilon}
\newcommand{\dd}{\,\mathrm{d}}
\begin{document}
\section*{UnsolvedMath ID: 30000525}
\textbf{Source problem code:} catalogue entry
\subsection{Definitions and mathematical statement}
Use the standard finite-valued continuous-integrand formulation. A function \(f:\R^{2\times2}\to\R\) is rank-one convex if
\[
f((1-t)A+tB)\leq(1-t)f(A)+tf(B)
\]
for every \(A,B\in\R^{2\times2}\) with \(\operatorname{rank}(A-B)\leq1\) and every \(t\in[0,1]\). Equivalently, it is convex along every line with direction \(a\otimes b\), \(a,b\in\R^2\).

Let \(Q=(0,1)^2\). Quasiconvexity in the standard smooth-test formulation means that, for every \(A\in\R^{2\times2}\) and every \(\varphi\in C_c^\infty(Q;\R^2)\),
\[
f(A)\leq\int_Q f(A+D\varphi(x))\,\dd x.
\]
The test perturbation vanishes near the boundary and \(|Q|=1\); on a general test domain the right side is averaged by its area. The usual continuous-integrand formulation may equivalently be stated with zero-boundary Lipschitz perturbations under the matching test conventions.

The source question is whether every rank-one convex integrand on the entire space \(\R^{2\times2}\) is quasiconvex in this sense. No determinant constraint, smaller matrix domain, quadratic-only restriction, or prescribed growth bound is added. The source does not explicitly state continuity or finite-valuedness; those standard conventions are made explicit here. Allowing extended-real integrands or imposing a restricted admissible-gradient class would require a separately specified question.

The supplied literature assessment checked 2026-08-21 says that the report identifies the exact case \(n=m=2\) as outstanding and that no later resolution of this case was located in that check. No result in a different matrix dimension is substituted for the two-by-two assertion.
\subsection{Short English statement}
Does convexity along every rank-one matrix direction force the variational quasiconvexity inequality for all two-dimensional vector perturbations? The target is the full two-by-two matrix space under the standard finite continuous-integrand convention.
\subsection{Sources}
\begin{itemize}
\item Local UnsolvedMath registry, permanent ID 30000525, catalogue entry, \emph{Rank-One Convexity Versus Quasiconvexity in Two Dimensions}.
\item \emph{Calculus of Variations} (2007), Oberwolfach Reports, pp. 1879--1940, DOI \url{https://doi.org/10.4171/owr/2006/31}.
\item Supplied statement-verification claim and literature triage checked 2026-08-21 cite only the official report; no separate exact-case theorem citation is supplied.
\end{itemize}
\end{document}
